% report-with-tables.tex — tabular material in LaTeX: simple tables, a table
% spanning columns, a multi-page table, and a table with merged cells.
% Compile with: pdflatex report-with-tables.tex   (run twice for references)
%
% The company, the quarters, and every number below are invented.

\documentclass[11pt,a4paper]{report}

\usepackage[utf8]{inputenc}
\usepackage[T1]{fontenc}
\usepackage{lmodern}
\usepackage{geometry}
\usepackage{booktabs}
\usepackage{longtable}
\usepackage{multirow}
\usepackage{array}
\usepackage{caption}

\geometry{margin=2.4cm}
\captionsetup{font=small}

\newcolumntype{R}{>{\raggedleft\arraybackslash}p{2.2cm}}

\title{Northwind Ferry Cooperative\\ Operating Report}
\author{Prepared by the Finance Working Group}
\date{Quarter ending 30 June 2027}

\begin{document}

\maketitle
\tableofcontents
\listoftables

\chapter{Summary}

The cooperative operated four routes across the quarter, carrying 214{,}880
passengers on 3{,}102 sailings. Two routes covered their direct costs; the
Halloway service did not, and the Kestrel Point service did so only after the
seasonal fare adjustment in May.

Every figure in this report is fabricated for demonstration purposes.

\chapter{Traffic}

\section{Passengers by route}

Table~\ref{tab:passengers} gives the passenger count for each route, split by
month.

\begin{table}[ht]
  \centering
  \caption{Passengers carried, by route and month}
  \label{tab:passengers}
  \begin{tabular}{l R R R R}
    \toprule
    Route          & April  & May    & June   & Total  \\
    \midrule
    Harbour Loop   & 24\,910 & 27\,640 & 31\,205 & 83\,755 \\
    Kestrel Point  & 18\,330 & 20\,115 & 24\,880 & 63\,325 \\
    Halloway       &  9\,470 &  9\,905 & 11\,240 & 30\,615 \\
    Night Crossing & 11\,020 & 12\,640 & 13\,525 & 37\,185 \\
    \midrule
    \textbf{Total} & \textbf{63\,730} & \textbf{70\,300} & \textbf{80\,850}
                   & \textbf{214\,880} \\
    \bottomrule
  \end{tabular}
\end{table}

\section{Load factor}

Table~\ref{tab:load} uses a merged first column to group the two daytime routes
against the two others.

\begin{table}[ht]
  \centering
  \caption{Average load factor by service class}
  \label{tab:load}
  \begin{tabular}{llrrr}
    \toprule
    Class & Route & April & May & June \\
    \midrule
    \multirow{2}{*}{Daytime}
          & Harbour Loop   & 61\,\% & 66\,\% & 74\,\% \\
          & Kestrel Point  & 54\,\% & 58\,\% & 69\,\% \\
    \midrule
    \multirow{2}{*}{Other}
          & Halloway       & 38\,\% & 39\,\% & 44\,\% \\
          & Night Crossing & 47\,\% & 52\,\% & 55\,\% \\
    \bottomrule
  \end{tabular}
\end{table}

\chapter{Finance}

\section{Direct cost recovery}

\begin{table}[ht]
  \centering
  \caption{Revenue against direct cost, quarter to 30 June (thousands)}
  \label{tab:recovery}
  \begin{tabular}{lrrrr}
    \toprule
    Route          & Revenue & Direct cost & Margin & Recovery \\
    \midrule
    Harbour Loop   & 1\,284 & 1\,011 & 273  & 127\,\% \\
    Kestrel Point  &   902  &   874  &  28  & 103\,\% \\
    Halloway       &   341  &   498  & -157 &  68\,\% \\
    Night Crossing &   556  &   540  &  16  & 103\,\% \\
    \midrule
    \textbf{Total} & \textbf{3\,083} & \textbf{2\,923} & \textbf{160}
                   & \textbf{105\,\%} \\
    \bottomrule
  \end{tabular}
\end{table}

\section{Sailings, cancellations, and causes}

The table below runs past a single page on purpose, so that a parser or a
converter meets a repeated header and a continuation footer.

\begin{longtable}{lllr}
  \caption{Cancelled sailings, quarter to 30 June}
  \label{tab:cancellations} \\
  \toprule
  Date & Route & Cause & Passengers affected \\
  \midrule
  \endfirsthead

  \multicolumn{4}{l}{\small\itshape Table \thetable\ continued from previous page} \\
  \toprule
  Date & Route & Cause & Passengers affected \\
  \midrule
  \endhead

  \midrule
  \multicolumn{4}{r}{\small\itshape continued on next page} \\
  \endfoot

  \bottomrule
  \endlastfoot

  02 Apr & Halloway       & Weather: swell        & 118 \\
  02 Apr & Night Crossing & Weather: swell        & 204 \\
  06 Apr & Kestrel Point  & Mechanical: gearbox   & 331 \\
  07 Apr & Kestrel Point  & Mechanical: gearbox   & 298 \\
  11 Apr & Harbour Loop   & Crew shortage         & 176 \\
  14 Apr & Halloway       & Weather: fog          &  84 \\
  15 Apr & Halloway       & Weather: fog          &  91 \\
  19 Apr & Night Crossing & Berth unavailable     & 143 \\
  23 Apr & Harbour Loop   & Weather: wind         & 260 \\
  28 Apr & Kestrel Point  & Crew shortage         & 187 \\
  03 May & Halloway       & Weather: swell        & 102 \\
  09 May & Night Crossing & Mechanical: generator & 221 \\
  10 May & Night Crossing & Mechanical: generator & 209 \\
  12 May & Harbour Loop   & Berth unavailable     & 314 \\
  17 May & Kestrel Point  & Weather: wind         & 246 \\
  21 May & Halloway       & Crew shortage         &  77 \\
  24 May & Harbour Loop   & Weather: wind         & 288 \\
  30 May & Night Crossing & Weather: fog          & 165 \\
  02 Jun & Kestrel Point  & Mechanical: steering  & 355 \\
  02 Jun & Halloway       & Mechanical: steering  &  96 \\
  08 Jun & Harbour Loop   & Crew shortage         & 201 \\
  13 Jun & Night Crossing & Weather: swell        & 188 \\
  16 Jun & Halloway       & Weather: fog          &  88 \\
  19 Jun & Kestrel Point  & Berth unavailable     & 271 \\
  22 Jun & Harbour Loop   & Weather: wind         & 305 \\
  25 Jun & Night Crossing & Crew shortage         & 174 \\
  27 Jun & Halloway       & Weather: swell        & 109 \\
  29 Jun & Kestrel Point  & Weather: wind         & 233 \\
\end{longtable}

\section{Cause summary}

Weather accounts for slightly more than half of all cancellations, and for the
majority of those on the two exposed routes. Mechanical causes cluster: the
Kestrel Point gearbox took two consecutive days in April, and the Night Crossing
generator took two in May.

\begin{table}[ht]
  \centering
  \caption{Cancellations by cause}
  \label{tab:causes}
  \begin{tabular}{lrr}
    \toprule
    Cause & Sailings & Share \\
    \midrule
    Weather            & 15 & 53.6\,\% \\
    Crew shortage      &  5 & 17.9\,\% \\
    Mechanical         &  5 & 17.9\,\% \\
    Berth unavailable  &  3 & 10.7\,\% \\
    \midrule
    \textbf{Total}     & \textbf{28} & \textbf{100.0\,\%} \\
    \bottomrule
  \end{tabular}
\end{table}

\chapter{Recommendations}

\begin{enumerate}
  \item Retire the Kestrel Point gearbox at the winter overhaul rather than
        waiting for the scheduled replacement in the following year.
  \item Review the Halloway timetable. Its recovery rate has not exceeded
        70\,\% in six quarters, and its cancellations are dominated by weather
        the timetable does not account for.
  \item Hold the seasonal fare adjustment for a full year before judging it.
        One quarter is not enough to separate it from the summer traffic.
\end{enumerate}

\end{document}
